\chapter{Método}\label{cap:metodo}

\section{Visión general}

El trabajo tiene dos tareas que comparten los datos y la validación (\cref{fig:flujo}). La primera, el rango esperado,
es un modelo residual pequeño que corre en cualquier CPU. La segunda, el pronóstico a 7 días, entrena redes neuronales
y corre solo en Kaggle. Las dos se evalúan con validación por sujetos: una persona que se usa para medir nunca entra a
ningún ajuste.

\begin{figure}[htbp]
\centering
\begin{tikzpicture}[
  caja/.style={draw=#1, line width=0.8pt, rounded corners=2pt, fill=white, align=center, font=\footnotesize,
    text width=3.05cm, minimum height=1.25cm, inner sep=3pt},
  caja/.default=tinta3,
  flecha/.style={-{Stealth[length=2.2mm]}, line width=0.8pt, color=tinta2},
  etiqueta/.style={font=\scriptsize\sffamily\color{tinta2}, align=center}]
  \node[caja=serieB] (reposo) at (0,1) {\textbf{Reposo en cama}\\ NASA NLSP/LSDA\\ \cifra{datos.sujetos} sujetos};
  \node[caja=serieB] (ls) at (0,-1) {\textbf{LifeSnaps}\\ reloj Fitbit\\ \cifra{datos.ls.personas} personas};
  \node[caja] (datos) at (3.95,0) {\texttt{nova\_tl.datos}\\ formato largo\\ 3 archivos parquet};
  \node[caja=serieA] (t1) at (7.9,1) {\textbf{Tarea 1}\\ rango esperado\\ ridge, en CPU};
  \node[caja=serieA] (t2) at (7.9,-1) {\textbf{Tarea 2}\\ pronóstico a 7 días\\ PatchTST, GPU T4};
  \node[caja] (sal) at (11.85,1) {\texttt{metricas.json}\\ tablas por sujeto\\ y gráficos};
  \node[caja] (hf) at (11.85,-1) {\textbf{Hugging Face}\\ \texttt{nova-tierra}\\ \texttt{nova-espacio},\\ con el corrector};
  \node[caja=azulnoche] (inf) at (11.85,3.2) {\textbf{Este informe}\\ \texttt{construir\_cifras.py}};
  \node[caja=tinta3, dashed] (nave) at (11.85,-3.2) {\textbf{NOVA a bordo}\\ servicio en el :8793\\ diseñado, no construido};
  \draw[flecha] (reposo.east) -- (datos.west);
  \draw[flecha] (ls.east) -- (datos.west);
  \draw[flecha] (datos.east) -- (t1.west);
  \draw[flecha] (datos.east) -- (t2.west);
  \draw[flecha] (t1.east) -- (sal.west);
  \draw[flecha] (t2.east) -- (sal.west);
  \draw[flecha] (t2.east) -- (hf.west);
  \draw[flecha] (sal.north) -- (inf.south);
  \draw[flecha, dashed] (hf.south) -- (nave.north);
  \begin{scope}[on background layer]
    \node[fill=fondo, rounded corners=4pt, inner sep=6pt, fit=(t1)(t2)(datos),
      label={[etiqueta, anchor=south]north:{kernel de Kaggle, validación por sujetos en 5 folds}}] {};
  \end{scope}
\end{tikzpicture}
\caption[Flujo del trabajo]{Flujo del trabajo. Los datos se llevan a un formato común; el kernel corre las dos tareas
con la misma validación por sujetos y deja métricas, tablas, gráficos y pesos. Este informe se arma desde esa salida.
La conexión con la nave está diseñada, pero no construida.}
\label{fig:flujo}
\end{figure}

\section{Rango esperado}\label{sec:rango}

El rango esperado responde a una pregunta simple: dado lo que se sabe de una persona, ¿qué valor se espera de ella hoy,
y cuánto se puede alejar de ese valor sin que sea raro? Es el modelo residual que el plan del proyecto ponía primero en
la lista de adaptaciones: una predicción de Tierra más un corrector pequeño y explicable.

\subsection{Línea base de Tierra}

Para cada sujeto $s$ y variable $v$, la línea base es la media y la desviación de sus días de pre-reposo:
\begin{equation}\label{eq:lineabase}
\mu_{s,v} = \frac{1}{n^{\text{pre}}_{s,v}} \sum_{d<0} x_{s,v,d}, \qquad
\sigma_{s,v} = \max\Bigl(\operatorname{de}_{d<0}(x_{s,v,d}),\; \tfrac{1}{2}\,\operatorname{mediana}_{s'\in\text{entr.}}\,
\operatorname{de}(x_{s',v})\Bigr).
\end{equation}
Se exigen al menos 3 días de pre-reposo. La desviación tiene un piso, la mitad de la mediana de las desviaciones de los
sujetos de entrenamiento, porque con 10 a 21 días una desviación casualmente chica dispararía alarmas. El
\emph{monitor de Tierra} es esa línea base tal cual: espera $\mu_{s,v}$ y acepta $\pm 1{,}96\,\sigma_{s,v}$.

\subsection{Corrector aprendido de otros sujetos}

En cama, el residuo respecto de la línea base, $r_{s,v,d} = x_{s,v,d} - \mu_{s,v}$, tiene una forma típica: el peso, por
ejemplo, baja rápido al principio y después más lento. El corrector poblacional la aprende con una regresión ridge
\parencite{hoerl1970} sobre dos rasgos del día en cama $d$, ajustada con los residuos de los sujetos de entrenamiento:
\begin{equation}\label{eq:corrector}
\Delta_v(d) = \beta_{0,v} + \beta_{1,v} \log(1+d) + \beta_{2,v}\, \frac{d}{30}, \qquad
\hat\beta_v = \arg\min_{\beta} \sum_{s,d} \bigl(r_{s,v,d} - \Delta_v(d)\bigr)^2 + \alpha\bigl(\beta_{1,v}^2 + \beta_{2,v}^2\bigr),
\end{equation}
con $\alpha = 1$. Son tres coeficientes por variable (\cref{tab:corrector}). El logaritmo captura el cambio rápido de
las primeras semanas y el término lineal la deriva lenta. Este corrector es lo único que se transfiere \emph{entre
personas}: lo que se aprendió del reposo de otros.

\subsection{Ajuste con los días propios}

Con los primeros $k$ días en cama de la propia persona se estima cuánto se aparta de lo típico. El ajuste personal es
la media del residuo que queda después del corrector, contraída hacia cero:
\begin{equation}\label{eq:personal}
\delta_{s,v}(k) = \frac{n}{n + \kappa}\cdot\frac{1}{n}\sum_{d<k}\bigl(r_{s,v,d} - \Delta_v(d)\bigr)
               = \frac{1}{n+\kappa}\sum_{d<k}\bigl(r_{s,v,d} - \Delta_v(d)\bigr),
\end{equation}
donde $n$ es el número de días medidos antes del día $k$ y $\kappa = 3$. La contracción equivale a sumar $\kappa$ días
ficticios con residuo cero: con pocos días el ajuste es tímido y con muchos se acerca a la media observada. El valor
esperado final es
\begin{equation}\label{eq:esperado}
\hat x_{s,v,d}(k) = \mu_{s,v} + \Delta_v(d) + \delta_{s,v}(k).
\end{equation}
Se prueban $k \in \{0, 3, 7, 14, 21, 28\}$. El modelo \enquote{adaptado} de los titulares es el de $k = 28$.

\subsection{Ancho del rango}

El monitor de Tierra usa $\pm1{,}96\,\sigma_{s,v}$, que supone que los días en cama se parecen a los del pre-reposo.
Los modelos adaptados usan un \emph{ancho calibrado}, igual para todas las personas de una variable: el cuantil 95 del
error absoluto fuera de muestra,
\begin{equation}\label{eq:ancho}
a_v = q_{0{,}95}\bigl\{\,|x_{s,v,d} - \hat x^{\text{(VC)}}_{s,v,d}|\; : \; s \in \text{entrenamiento},\; d \ge 28 \,\bigr\},
\end{equation}
donde $\hat x^{\text{(VC)}}$ es la predicción para cada sujeto de entrenamiento hecha con un corrector que no lo vio,
por validación cruzada en 4 grupos de sujetos. Es la misma idea de la predicción conforme partida
\parencite{vovk2005,lei2018}: calibrar el ancho con residuos de datos que el modelo no usó. Como en la predicción
conforme, la garantía es \emph{marginal}: el 5\pct{} de días fuera se cumple en promedio sobre la población, no en
cada persona.

\subsection{Regla de alarma}

Un día queda \emph{fuera} si $|x_{s,v,d} - \hat x_{s,v,d}| > a$. Como un día fuera puede ser ruido, la alarma se
levanta cuando al menos 2 de 3 días seguidos de calendario quedan fuera:
\begin{equation}\label{eq:regla}
\text{alarma}_{s,v,d} = \mathbb{1}\Bigl[\,\textstyle\sum_{j=d-2}^{d} \text{fuera}_{s,v,j} \ge 2\Bigr],
\end{equation}
definida solo si los tres días se midieron. Es una regla de la familia de las cartas de control: se parece a una de las
pruebas de \textcite{nelson1984}, \enquote{2 de 3 puntos seguidos más allá de 2 sigmas}, aunque aquí no se exige que
estén del mismo lado. Para cambios sostenidos y pequeños existen reglas acumulativas como la CUSUM de
\textcite{page1954}, que este trabajo no prueba.

\subsection{Seis modelos para separar el centro del ancho}

Si el modelo adaptado tuviera menos falsas alarmas solo por usar un ancho calibrado, el mérito sería de la calibración y
no de la adaptación. Para separar las dos cosas se cruzan centros y anchos (\cref{tab:modelos-rango}).

\begin{table}[htbp]
\caption[Modelos del rango esperado]{Modelos del rango esperado. Cruzar centro y ancho permite ver cuánto de la mejora
viene de mover el centro y cuánto de calibrar el ancho.}
\label{tab:modelos-rango}
\small
\begin{tabularx}{\textwidth}{@{}l >{\raggedright\arraybackslash}X l l@{}}
\toprule
Modelo & Centro & $k$ & Ancho \\
\midrule
\modelo{tierra} & $\mu_{s,v}$ & 0 & $\pm1{,}96\,\sigma_{s,v}$ \\
\modelo{tierra (ancho calibrado)} & $\mu_{s,v}$ & 0 & calibrado \\
\modelo{tierra+personal} & $\mu_{s,v} + \delta_{s,v}(k)$, sin corrector & 3 a 28 & calibrado \\
\modelo{tierra+poblacion} & $\mu_{s,v} + \Delta_v(d)$ & 0 & calibrado \\
\modelo{tierra+poblacion+personal} & $\mu_{s,v} + \Delta_v(d) + \delta_{s,v}(k)$ & 3 a 28 & calibrado \\
\modelo{adaptado (ancho de Tierra)} & $\mu_{s,v} + \Delta_v(d) + \delta_{s,v}(28)$ & 28 & $\pm1{,}96\,\sigma_{s,v}$ \\
\bottomrule
\end{tabularx}
\end{table}

\subsection{Escalones sintéticos para medir la detección}\label{sec:escalones}

En el reposo en cama no hay anomalías marcadas, así que para medir si un modelo detecta algo hay que ponerlo. En cada
sujeto de prueba y variable se inyectan hasta $\min(3, \lfloor n/12 \rfloor)$ escalones, donde $n$ es el número de días
evaluados: bloques de 3 días seguidos medidos, desde el día 28, que no se tocan entre sí y dejan al menos un día libre
entre uno y otro. Cada escalón suma o resta, con signo al azar, la magnitud de la \cref{tab:variables}: 10\un{lpm} de
pulso, 15 y 10\un{mmHg} de presión sistólica y diastólica, 4\un{rpm} de respiración, 0,8\un{°C} de temperatura,
1,5\un{kg} de peso y 1000\un{g/día} de diuresis o agua. Solo la del pulso tiene fuente, los umbrales de NOVA
(\enquote{pulso en reposo +10\un{lpm} en 3 días}); las demás son supuestos, elegidas como un cambio clínicamente
notable.

Un escalón cuenta como detectado si la regla de 2 de 3 se activa en su último día. Como el signo es al azar, algunos
escalones van a favor de la deriva del reposo (por ejemplo, el peso que baja aún más) y otros en contra; la detección
se informa también por separado, porque un monitor desplazado detecta con facilidad lo que va a favor y casi nada de lo
que va en contra.

\subsection{Métricas y placebo}

Todo se evalúa en los días desde el 28 en cama, para que todos los $k$ se comparen en el mismo tramo:
\begin{description}
\item[Falsas alarmas diarias] porcentaje de días normales (sin escalón) que quedan fuera del rango.
\item[Falsas alarmas con la regla] porcentaje de bloques de 3 días normales en que la regla de 2 de 3 se activa.
\item[Detección] porcentaje de escalones detectados, en total, a favor y en contra de la deriva.
\item[Ancho] ancho medio del rango, en las unidades de la variable.
\item[Error] error absoluto medio del valor esperado, en días normales.
\end{description}
Todas son cocientes de sumas sobre sujetos, lo que permite remuestrear personas (\cref{sec:validacion}).

El \emph{placebo} mide cuántas falsas alarmas da la regla de Tierra sin ningún efecto del reposo: dentro del pre-reposo,
cada día se compara con la media y la desviación de los \emph{demás} días del pre-reposo (con el mismo piso), y cada
bloque de 3 días con los días restantes. Es el piso de falsas alarmas que tiene esa regla solo por estimar la
variación de una persona con 10 a 21 días.

\section{Pronóstico a 7 días}\label{sec:pronostico}

\subsection{Ventanas}

Cada serie diaria de una persona y una variable se corta en ventanas: 14 días de contexto y los 7 siguientes como
objetivo, una ventana por día de inicio. Los días sin medición del contexto se rellenan con el último valor observado
(o el siguiente, al inicio) y quedan marcados como no observados; en el objetivo quedan con máscara cero, así la
pérdida y las métricas solo cuentan lo medido. El objetivo no puede cruzar de una fase a otra. Se exigen al menos 5
días observados en el contexto en el reposo y 7 en LifeSnaps. Como el pre-reposo dura 10 a 21 días y no alcanza para
14 de contexto más 7 de objetivo, se permiten ventanas cerca del inicio de la serie con el contexto completado a la
izquierda con días no observados. Las ventanas que se evalúan en cama empiezan su objetivo desde el día 14, para que el
contexto entero esté en cama.

\subsection{PatchTST}

El modelo es \texttt{PatchTSTForPrediction} de \texttt{transformers} \parencite{nie2023,wolf2020} con un canal: es
univariado y compartido entre variables. La \cref{fig:patchtst} muestra su arquitectura y la \cref{tab:patchtst} su
configuración.

\begin{figure}[htbp]
\centering
\begin{tikzpicture}[
  bloque/.style={draw=#1, line width=0.8pt, rounded corners=2pt, fill=white, align=center, font=\footnotesize,
    minimum height=1.55cm, inner sep=3pt},
  flecha/.style={-{Stealth[length=2mm]}, line width=0.8pt, color=tinta2}]
  % serie de entrada
  \begin{scope}[shift={(0,0)}]
    \foreach \i/\y in {0/0.30,1/0.42,2/0.36,3/0.55,4/0.50,5/0.62,6/0.48,7/0.58,8/0.70,9/0.64,10/0.72,11/0.60,12/0.68,13/0.75}{
      \fill[tinta2] (\i*0.12, \y) circle[radius=0.9pt];
    }
    \draw[tinta3, line width=0.5pt] (0,0.30) -- (0.12,0.42) -- (0.24,0.36) -- (0.36,0.55) -- (0.48,0.50) -- (0.60,0.62)
      -- (0.72,0.48) -- (0.84,0.58) -- (0.96,0.70) -- (1.08,0.64) -- (1.20,0.72) -- (1.32,0.60) -- (1.44,0.68) -- (1.56,0.75);
    \node[font=\scriptsize\sffamily\color{tinta2}, align=center] at (0.78,-0.25) {14 días};
  \end{scope}
  \node[bloque=tinta3, text width=1.65cm] (esc) at (3.0,0.35) {escala\\ \texttt{std}};
  \node[bloque=tinta3, text width=1.9cm] (par) at (5.35,0.35) {\cifra{modelo.parches} parches\\ de \cifra{cfg.patch.length},
    paso \cifra{cfg.patch.stride}};
  \node[bloque=tinta3, text width=2.0cm] (emb) at (7.85,0.35) {proyección a\\ $d=\cifra{cfg.d.model}$ +\\ posición sincos};
  \node[bloque=serieA, text width=2.35cm, minimum height=2.1cm] (enc) at (10.65,0.35) {\cifra{cfg.num.hidden.layers}
    capas:\\ atención de \cifra{cfg.num.attention.heads} cabezas,\\ FFN de \cifra{cfg.ffn.dim},\\ BatchNorm};
  \node[bloque=serieB, text width=1.75cm] (cab) at (13.35,0.35) {promedio\\ de parches,\\ lineal a 7};
  \draw[flecha] (1.75,0.5) -- (esc.west);
  \draw[flecha] (esc) -- (par);
  \draw[flecha] (par) -- (emb);
  \draw[flecha] (emb) -- (enc);
  \draw[flecha] (enc) -- (cab);
  \draw[flecha] (cab.south) -- ++(0,-0.75) node[below, font=\scriptsize\sffamily\color{tinta2}] {7 días};
  \draw[decorate, decoration={brace, amplitude=5pt, mirror}, tinta2] (6.3,-0.75) -- (11.85,-0.75)
    node[midway, below=6pt, font=\scriptsize\sffamily\color{tinta2}, align=center] {backbone: congelado en
    \enquote{solo la cabeza},\\ salvo las estadísticas de BatchNorm};
  \draw[decorate, decoration={brace, amplitude=5pt}, serieB] (12.45,1.45) -- (14.25,1.45)
    node[midway, above=6pt, font=\scriptsize\sffamily\color{tinta2}] {cabeza: \cifra{modelo.cabeza} parámetros};
\end{tikzpicture}
\caption[Arquitectura del PatchTST]{Arquitectura del PatchTST de NOVA. La serie se normaliza, se corta en parches que
funcionan como fichas, pasa por un codificador transformer y una cabeza lineal produce los 7 días. Tiene
\cifra{modelo.parametros} parámetros entrenables, de los que \cifra{modelo.cabeza} (\cifra{modelo.cabeza.pct}\pct) son
de la cabeza.}
\label{fig:patchtst}
\end{figure}

\begin{table}[htbp]
\caption[Configuración del PatchTST]{Configuración del PatchTST, leída de \texttt{config.json} de los pesos
publicados.}
\label{tab:patchtst}
\small
\begin{tabular}{@{}ll@{\hspace{2.2em}}ll@{}}
\toprule
Parámetro & Valor & Parámetro & Valor \\
\midrule
Contexto & \cifra{cfg.context.length} días & Dimensión del modelo & \cifra{cfg.d.model} \\
Horizonte & \cifra{cfg.prediction.length} días & Cabezas de atención & \cifra{cfg.num.attention.heads} \\
Largo del parche & \cifra{cfg.patch.length} días & Capas del codificador & \cifra{cfg.num.hidden.layers} \\
Paso del parche & \cifra{cfg.patch.stride} días & Dimensión de la FFN & \cifra{cfg.ffn.dim} \\
Parches por ventana & \cifra{modelo.parches} & Dropout / de la cabeza & \cifra{cfg.dropout} / \cifra{cfg.head.dropout} \\
Escalado & \texttt{std} & Normalización & BatchNorm, antes de cada subcapa \\
Posición & sinusoidal & Agregación & promedio de parches \\
Activación & GELU & Parámetros entrenables & \cifra{modelo.parametros} \\
Canales & 1 (compartido) & Buffers de BatchNorm & \cifra{modelo.buffers} \\
\bottomrule
\end{tabular}
\end{table}

La pérdida que trae \texttt{PatchTSTForPrediction} es un MSE sin máscara, así que se calcula aparte: el error al
cuadrado, dividido por el ruido diario típico de cada variable $s_v$ (\cref{tab:variables}), solo en los días
observados del objetivo,
\begin{equation}\label{eq:perdida}
\mathcal{L} = \frac{\sum_{i,h} m_{i,h}\,\bigl(\hat y_{i,h} - y_{i,h}\bigr)^2 / s_{v(i)}^2}{\sum_{i,h} m_{i,h}} ,
\end{equation}
donde $m_{i,h}$ vale 1 si el día $h$ de la ventana $i$ se midió. Dividir por $s_v$ hace que el pulso y la diuresis
pesen parecido aunque sus unidades sean muy distintas. El ruido se calcula en cada fold solo con LifeSnaps y el
pre-reposo de los sujetos de entrenamiento.

El entrenamiento usa AdamW \parencite{kingma2015,loshchilov2019} con decaimiento de pesos de $10^{-4}$, lotes de 512
ventanas y recorte del gradiente a norma 1. Se detiene cuando la pérdida de validación no mejora en 6 épocas
\parencite{prechelt1998}, y se queda con el mejor estado. El dropout \parencite{srivastava2014} está activo en
entrenamiento.

\subsection{Brazos del experimento}

La \cref{tab:brazos} define los modelos que se comparan. El modelo de Tierra se entrena de nuevo en cada fold, sin el
pre-reposo de los sujetos de prueba. Los dos brazos de adaptación parten de ese modelo y se ajustan con los días en
cama de los sujetos de entrenamiento. El sujeto evaluado nunca aporta sus propios días: el pronóstico se adapta al
reposo en general, no a cada persona.

\begin{table}[htbp]
\caption[Brazos del pronóstico]{Brazos del pronóstico. \enquote{Entr.}: sujetos de entrenamiento del fold, sin los 4
de validación. Los dos brazos de adaptación y los dos desde cero no se entrenaron en las mismas condiciones; la
\cref{sec:h1} discute por qué importa.}
\label{tab:brazos}
\small
\begin{tabularx}{\textwidth}{@{}>{\raggedright\arraybackslash}p{3.3cm} >{\raggedright\arraybackslash}X
  >{\raggedright\arraybackslash}p{2.2cm} l r@{}}
\toprule
Brazo & Datos de entrenamiento & Parte de & lr & Tope \\
\midrule
Tierra (base) & LifeSnaps (\cifra{datos.ls.entrenan} personas) y pre-reposo de entr. & azar & \texttt{1e-3} & 60 \\
Tierra sin adaptar & ninguno más & Tierra & -- & -- \\
Tierra → espacio, solo la cabeza & días en cama de entr.; backbone sin gradientes, modelo en modo entrenamiento & Tierra
& \texttt{1e-3} & 30 \\
Tierra → espacio, ajuste completo & días en cama de entr.; todos los parámetros & Tierra & \texttt{2e-4} & 30 \\
Reposo, desde cero & días en cama de entr. & azar & \texttt{1e-3} & 60 \\
Pre-reposo + reposo, desde cero & pre-reposo y días en cama de entr., sin LifeSnaps & azar & \texttt{1e-3} & 60 \\
Chronos-2 sin adaptar & ninguno; contexto con los huecos rellenados & \texttt{amazon/\allowbreak chronos-2} & -- & -- \\
Persistencia & repite el último día observado & -- & -- & -- \\
Media de 7 días & repite la media de los últimos 7 días & -- & -- & -- \\
\bottomrule
\end{tabularx}
\end{table}

La métrica es el MAE de cada modelo dividido por el MAE de la persistencia, por variable, promediado entre las siete
variables sin el agua \parencite{hyndman2006}. Bajo 1, el modelo se equivoca menos que repetir el último día. Las
comparaciones entre modelos son diferencias pareadas por sujeto de ese promedio.

\section{Validación}\label{sec:validacion}

\subsection{Por sujetos, estratificada por campaña}

Los días de una misma persona se parecen entre sí, así que validar con días mezclados sobreestima el desempeño
\parencite{saeb2017}. Aquí se valida por sujetos: dentro de cada campaña, los sujetos se permutan con una semilla fija
y se reparten por turno en 5 folds, de modo que cada fold tiene sujetos de las tres campañas en proporción
(\cref{fig:folds}).

\begin{figure}[htbp]
\centering
\begin{tikzpicture}[x=4.6mm, y=7mm]
\node[anchor=east, font=\footnotesize\sffamily\color{tinta2}] at (-0.4,0) {fold 0};
\draw[serieB!45, line width=0.6pt, rounded corners=0.6pt] (0.04,-0.28) rectangle (0.76,0.28);
\draw[serieB!45, line width=0.6pt, rounded corners=0.6pt] (1.04,-0.28) rectangle (1.76,0.28);
\fill[serieB, rounded corners=0.6pt] (2,-0.32) rectangle (2.8,0.32);
\draw[serieC!45, line width=0.6pt, rounded corners=0.6pt] (3.04,-0.28) rectangle (3.76,0.28);
\draw[serieC!45, line width=0.6pt, rounded corners=0.6pt] (4.04,-0.28) rectangle (4.76,0.28);
\draw[serieC!45, line width=0.6pt, rounded corners=0.6pt] (5.04,-0.28) rectangle (5.76,0.28);
\draw[serieA!45, line width=0.6pt, rounded corners=0.6pt] (6.04,-0.28) rectangle (6.76,0.28);
\draw[serieA!45, line width=0.6pt, rounded corners=0.6pt] (7.04,-0.28) rectangle (7.76,0.28);
\draw[serieA!45, line width=0.6pt, rounded corners=0.6pt] (8.04,-0.28) rectangle (8.76,0.28);
\draw[serieA!45, line width=0.6pt, rounded corners=0.6pt] (9.04,-0.28) rectangle (9.76,0.28);
\draw[serieA!45, line width=0.6pt, rounded corners=0.6pt] (10.04,-0.28) rectangle (10.76,0.28);
\draw[serieA!45, line width=0.6pt, rounded corners=0.6pt] (11.04,-0.28) rectangle (11.76,0.28);
\fill[serieA, rounded corners=0.6pt] (12,-0.32) rectangle (12.8,0.32);
\draw[serieA!45, line width=0.6pt, rounded corners=0.6pt] (13.04,-0.28) rectangle (13.76,0.28);
\fill[serieA, rounded corners=0.6pt] (14,-0.32) rectangle (14.8,0.32);
\fill[serieA, rounded corners=0.6pt] (15,-0.32) rectangle (15.8,0.32);
\fill[serieA, rounded corners=0.6pt] (16,-0.32) rectangle (16.8,0.32);
\draw[serieA!45, line width=0.6pt, rounded corners=0.6pt] (17.04,-0.28) rectangle (17.76,0.28);
\draw[serieA!45, line width=0.6pt, rounded corners=0.6pt] (18.04,-0.28) rectangle (18.76,0.28);
\fill[serieA, rounded corners=0.6pt] (19,-0.32) rectangle (19.8,0.32);
\draw[serieA!45, line width=0.6pt, rounded corners=0.6pt] (20.04,-0.28) rectangle (20.76,0.28);
\draw[serieA!45, line width=0.6pt, rounded corners=0.6pt] (21.04,-0.28) rectangle (21.76,0.28);
\draw[serieA!45, line width=0.6pt, rounded corners=0.6pt] (22.04,-0.28) rectangle (22.76,0.28);
\draw[serieA!45, line width=0.6pt, rounded corners=0.6pt] (23.04,-0.28) rectangle (23.76,0.28);
\draw[serieA!45, line width=0.6pt, rounded corners=0.6pt] (24.04,-0.28) rectangle (24.76,0.28);
\draw[serieA!45, line width=0.6pt, rounded corners=0.6pt] (25.04,-0.28) rectangle (25.76,0.28);
\draw[serieA!45, line width=0.6pt, rounded corners=0.6pt] (26.04,-0.28) rectangle (26.76,0.28);
\draw[serieA!45, line width=0.6pt, rounded corners=0.6pt] (27.04,-0.28) rectangle (27.76,0.28);
\draw[serieA!45, line width=0.6pt, rounded corners=0.6pt] (28.04,-0.28) rectangle (28.76,0.28);
\node[anchor=east, font=\footnotesize\sffamily\color{tinta2}] at (-0.4,-1) {fold 1};
\fill[serieB, rounded corners=0.6pt] (0,-1.32) rectangle (0.8,-0.6799999999999999);
\draw[serieB!45, line width=0.6pt, rounded corners=0.6pt] (1.04,-1.28) rectangle (1.76,-0.72);
\draw[serieB!45, line width=0.6pt, rounded corners=0.6pt] (2.04,-1.28) rectangle (2.76,-0.72);
\draw[serieC!45, line width=0.6pt, rounded corners=0.6pt] (3.04,-1.28) rectangle (3.76,-0.72);
\draw[serieC!45, line width=0.6pt, rounded corners=0.6pt] (4.04,-1.28) rectangle (4.76,-0.72);
\fill[serieC, rounded corners=0.6pt] (5,-1.32) rectangle (5.8,-0.6799999999999999);
\fill[serieA, rounded corners=0.6pt] (6,-1.32) rectangle (6.8,-0.6799999999999999);
\draw[serieA!45, line width=0.6pt, rounded corners=0.6pt] (7.04,-1.28) rectangle (7.76,-0.72);
\draw[serieA!45, line width=0.6pt, rounded corners=0.6pt] (8.04,-1.28) rectangle (8.76,-0.72);
\draw[serieA!45, line width=0.6pt, rounded corners=0.6pt] (9.04,-1.28) rectangle (9.76,-0.72);
\draw[serieA!45, line width=0.6pt, rounded corners=0.6pt] (10.04,-1.28) rectangle (10.76,-0.72);
\draw[serieA!45, line width=0.6pt, rounded corners=0.6pt] (11.04,-1.28) rectangle (11.76,-0.72);
\draw[serieA!45, line width=0.6pt, rounded corners=0.6pt] (12.04,-1.28) rectangle (12.76,-0.72);
\fill[serieA, rounded corners=0.6pt] (13,-1.32) rectangle (13.8,-0.6799999999999999);
\draw[serieA!45, line width=0.6pt, rounded corners=0.6pt] (14.04,-1.28) rectangle (14.76,-0.72);
\draw[serieA!45, line width=0.6pt, rounded corners=0.6pt] (15.04,-1.28) rectangle (15.76,-0.72);
\draw[serieA!45, line width=0.6pt, rounded corners=0.6pt] (16.04,-1.28) rectangle (16.76,-0.72);
\fill[serieA, rounded corners=0.6pt] (17,-1.32) rectangle (17.8,-0.6799999999999999);
\draw[serieA!45, line width=0.6pt, rounded corners=0.6pt] (18.04,-1.28) rectangle (18.76,-0.72);
\draw[serieA!45, line width=0.6pt, rounded corners=0.6pt] (19.04,-1.28) rectangle (19.76,-0.72);
\draw[serieA!45, line width=0.6pt, rounded corners=0.6pt] (20.04,-1.28) rectangle (20.76,-0.72);
\draw[serieA!45, line width=0.6pt, rounded corners=0.6pt] (21.04,-1.28) rectangle (21.76,-0.72);
\fill[serieA, rounded corners=0.6pt] (22,-1.32) rectangle (22.8,-0.6799999999999999);
\draw[serieA!45, line width=0.6pt, rounded corners=0.6pt] (23.04,-1.28) rectangle (23.76,-0.72);
\draw[serieA!45, line width=0.6pt, rounded corners=0.6pt] (24.04,-1.28) rectangle (24.76,-0.72);
\draw[serieA!45, line width=0.6pt, rounded corners=0.6pt] (25.04,-1.28) rectangle (25.76,-0.72);
\draw[serieA!45, line width=0.6pt, rounded corners=0.6pt] (26.04,-1.28) rectangle (26.76,-0.72);
\draw[serieA!45, line width=0.6pt, rounded corners=0.6pt] (27.04,-1.28) rectangle (27.76,-0.72);
\draw[serieA!45, line width=0.6pt, rounded corners=0.6pt] (28.04,-1.28) rectangle (28.76,-0.72);
\node[anchor=east, font=\footnotesize\sffamily\color{tinta2}] at (-0.4,-2) {fold 2};
\draw[serieB!45, line width=0.6pt, rounded corners=0.6pt] (0.04,-2.2800000000000002) rectangle (0.76,-1.72);
\fill[serieB, rounded corners=0.6pt] (1,-2.32) rectangle (1.8,-1.68);
\draw[serieB!45, line width=0.6pt, rounded corners=0.6pt] (2.04,-2.2800000000000002) rectangle (2.76,-1.72);
\draw[serieC!45, line width=0.6pt, rounded corners=0.6pt] (3.04,-2.2800000000000002) rectangle (3.76,-1.72);
\fill[serieC, rounded corners=0.6pt] (4,-2.32) rectangle (4.8,-1.68);
\draw[serieC!45, line width=0.6pt, rounded corners=0.6pt] (5.04,-2.2800000000000002) rectangle (5.76,-1.72);
\draw[serieA!45, line width=0.6pt, rounded corners=0.6pt] (6.04,-2.2800000000000002) rectangle (6.76,-1.72);
\draw[serieA!45, line width=0.6pt, rounded corners=0.6pt] (7.04,-2.2800000000000002) rectangle (7.76,-1.72);
\draw[serieA!45, line width=0.6pt, rounded corners=0.6pt] (8.04,-2.2800000000000002) rectangle (8.76,-1.72);
\draw[serieA!45, line width=0.6pt, rounded corners=0.6pt] (9.04,-2.2800000000000002) rectangle (9.76,-1.72);
\draw[serieA!45, line width=0.6pt, rounded corners=0.6pt] (10.04,-2.2800000000000002) rectangle (10.76,-1.72);
\fill[serieA, rounded corners=0.6pt] (11,-2.32) rectangle (11.8,-1.68);
\draw[serieA!45, line width=0.6pt, rounded corners=0.6pt] (12.04,-2.2800000000000002) rectangle (12.76,-1.72);
\draw[serieA!45, line width=0.6pt, rounded corners=0.6pt] (13.04,-2.2800000000000002) rectangle (13.76,-1.72);
\draw[serieA!45, line width=0.6pt, rounded corners=0.6pt] (14.04,-2.2800000000000002) rectangle (14.76,-1.72);
\draw[serieA!45, line width=0.6pt, rounded corners=0.6pt] (15.04,-2.2800000000000002) rectangle (15.76,-1.72);
\draw[serieA!45, line width=0.6pt, rounded corners=0.6pt] (16.04,-2.2800000000000002) rectangle (16.76,-1.72);
\draw[serieA!45, line width=0.6pt, rounded corners=0.6pt] (17.04,-2.2800000000000002) rectangle (17.76,-1.72);
\draw[serieA!45, line width=0.6pt, rounded corners=0.6pt] (18.04,-2.2800000000000002) rectangle (18.76,-1.72);
\draw[serieA!45, line width=0.6pt, rounded corners=0.6pt] (19.04,-2.2800000000000002) rectangle (19.76,-1.72);
\draw[serieA!45, line width=0.6pt, rounded corners=0.6pt] (20.04,-2.2800000000000002) rectangle (20.76,-1.72);
\fill[serieA, rounded corners=0.6pt] (21,-2.32) rectangle (21.8,-1.68);
\draw[serieA!45, line width=0.6pt, rounded corners=0.6pt] (22.04,-2.2800000000000002) rectangle (22.76,-1.72);
\draw[serieA!45, line width=0.6pt, rounded corners=0.6pt] (23.04,-2.2800000000000002) rectangle (23.76,-1.72);
\draw[serieA!45, line width=0.6pt, rounded corners=0.6pt] (24.04,-2.2800000000000002) rectangle (24.76,-1.72);
\draw[serieA!45, line width=0.6pt, rounded corners=0.6pt] (25.04,-2.2800000000000002) rectangle (25.76,-1.72);
\fill[serieA, rounded corners=0.6pt] (26,-2.32) rectangle (26.8,-1.68);
\draw[serieA!45, line width=0.6pt, rounded corners=0.6pt] (27.04,-2.2800000000000002) rectangle (27.76,-1.72);
\fill[serieA, rounded corners=0.6pt] (28,-2.32) rectangle (28.8,-1.68);
\node[anchor=east, font=\footnotesize\sffamily\color{tinta2}] at (-0.4,-3) {fold 3};
\draw[serieB!45, line width=0.6pt, rounded corners=0.6pt] (0.04,-3.2800000000000002) rectangle (0.76,-2.7199999999999998);
\draw[serieB!45, line width=0.6pt, rounded corners=0.6pt] (1.04,-3.2800000000000002) rectangle (1.76,-2.7199999999999998);
\draw[serieB!45, line width=0.6pt, rounded corners=0.6pt] (2.04,-3.2800000000000002) rectangle (2.76,-2.7199999999999998);
\fill[serieC, rounded corners=0.6pt] (3,-3.32) rectangle (3.8,-2.68);
\draw[serieC!45, line width=0.6pt, rounded corners=0.6pt] (4.04,-3.2800000000000002) rectangle (4.76,-2.7199999999999998);
\draw[serieC!45, line width=0.6pt, rounded corners=0.6pt] (5.04,-3.2800000000000002) rectangle (5.76,-2.7199999999999998);
\draw[serieA!45, line width=0.6pt, rounded corners=0.6pt] (6.04,-3.2800000000000002) rectangle (6.76,-2.7199999999999998);
\draw[serieA!45, line width=0.6pt, rounded corners=0.6pt] (7.04,-3.2800000000000002) rectangle (7.76,-2.7199999999999998);
\fill[serieA, rounded corners=0.6pt] (8,-3.32) rectangle (8.8,-2.68);
\fill[serieA, rounded corners=0.6pt] (9,-3.32) rectangle (9.8,-2.68);
\draw[serieA!45, line width=0.6pt, rounded corners=0.6pt] (10.04,-3.2800000000000002) rectangle (10.76,-2.7199999999999998);
\draw[serieA!45, line width=0.6pt, rounded corners=0.6pt] (11.04,-3.2800000000000002) rectangle (11.76,-2.7199999999999998);
\draw[serieA!45, line width=0.6pt, rounded corners=0.6pt] (12.04,-3.2800000000000002) rectangle (12.76,-2.7199999999999998);
\draw[serieA!45, line width=0.6pt, rounded corners=0.6pt] (13.04,-3.2800000000000002) rectangle (13.76,-2.7199999999999998);
\draw[serieA!45, line width=0.6pt, rounded corners=0.6pt] (14.04,-3.2800000000000002) rectangle (14.76,-2.7199999999999998);
\draw[serieA!45, line width=0.6pt, rounded corners=0.6pt] (15.04,-3.2800000000000002) rectangle (15.76,-2.7199999999999998);
\draw[serieA!45, line width=0.6pt, rounded corners=0.6pt] (16.04,-3.2800000000000002) rectangle (16.76,-2.7199999999999998);
\draw[serieA!45, line width=0.6pt, rounded corners=0.6pt] (17.04,-3.2800000000000002) rectangle (17.76,-2.7199999999999998);
\fill[serieA, rounded corners=0.6pt] (18,-3.32) rectangle (18.8,-2.68);
\draw[serieA!45, line width=0.6pt, rounded corners=0.6pt] (19.04,-3.2800000000000002) rectangle (19.76,-2.7199999999999998);
\draw[serieA!45, line width=0.6pt, rounded corners=0.6pt] (20.04,-3.2800000000000002) rectangle (20.76,-2.7199999999999998);
\draw[serieA!45, line width=0.6pt, rounded corners=0.6pt] (21.04,-3.2800000000000002) rectangle (21.76,-2.7199999999999998);
\draw[serieA!45, line width=0.6pt, rounded corners=0.6pt] (22.04,-3.2800000000000002) rectangle (22.76,-2.7199999999999998);
\fill[serieA, rounded corners=0.6pt] (23,-3.32) rectangle (23.8,-2.68);
\draw[serieA!45, line width=0.6pt, rounded corners=0.6pt] (24.04,-3.2800000000000002) rectangle (24.76,-2.7199999999999998);
\fill[serieA, rounded corners=0.6pt] (25,-3.32) rectangle (25.8,-2.68);
\draw[serieA!45, line width=0.6pt, rounded corners=0.6pt] (26.04,-3.2800000000000002) rectangle (26.76,-2.7199999999999998);
\draw[serieA!45, line width=0.6pt, rounded corners=0.6pt] (27.04,-3.2800000000000002) rectangle (27.76,-2.7199999999999998);
\draw[serieA!45, line width=0.6pt, rounded corners=0.6pt] (28.04,-3.2800000000000002) rectangle (28.76,-2.7199999999999998);
\node[anchor=east, font=\footnotesize\sffamily\color{tinta2}] at (-0.4,-4) {fold 4};
\draw[serieB!45, line width=0.6pt, rounded corners=0.6pt] (0.04,-4.28) rectangle (0.76,-3.7199999999999998);
\draw[serieB!45, line width=0.6pt, rounded corners=0.6pt] (1.04,-4.28) rectangle (1.76,-3.7199999999999998);
\draw[serieB!45, line width=0.6pt, rounded corners=0.6pt] (2.04,-4.28) rectangle (2.76,-3.7199999999999998);
\draw[serieC!45, line width=0.6pt, rounded corners=0.6pt] (3.04,-4.28) rectangle (3.76,-3.7199999999999998);
\draw[serieC!45, line width=0.6pt, rounded corners=0.6pt] (4.04,-4.28) rectangle (4.76,-3.7199999999999998);
\draw[serieC!45, line width=0.6pt, rounded corners=0.6pt] (5.04,-4.28) rectangle (5.76,-3.7199999999999998);
\draw[serieA!45, line width=0.6pt, rounded corners=0.6pt] (6.04,-4.28) rectangle (6.76,-3.7199999999999998);
\fill[serieA, rounded corners=0.6pt] (7,-4.32) rectangle (7.8,-3.68);
\draw[serieA!45, line width=0.6pt, rounded corners=0.6pt] (8.04,-4.28) rectangle (8.76,-3.7199999999999998);
\draw[serieA!45, line width=0.6pt, rounded corners=0.6pt] (9.04,-4.28) rectangle (9.76,-3.7199999999999998);
\fill[serieA, rounded corners=0.6pt] (10,-4.32) rectangle (10.8,-3.68);
\draw[serieA!45, line width=0.6pt, rounded corners=0.6pt] (11.04,-4.28) rectangle (11.76,-3.7199999999999998);
\draw[serieA!45, line width=0.6pt, rounded corners=0.6pt] (12.04,-4.28) rectangle (12.76,-3.7199999999999998);
\draw[serieA!45, line width=0.6pt, rounded corners=0.6pt] (13.04,-4.28) rectangle (13.76,-3.7199999999999998);
\draw[serieA!45, line width=0.6pt, rounded corners=0.6pt] (14.04,-4.28) rectangle (14.76,-3.7199999999999998);
\draw[serieA!45, line width=0.6pt, rounded corners=0.6pt] (15.04,-4.28) rectangle (15.76,-3.7199999999999998);
\draw[serieA!45, line width=0.6pt, rounded corners=0.6pt] (16.04,-4.28) rectangle (16.76,-3.7199999999999998);
\draw[serieA!45, line width=0.6pt, rounded corners=0.6pt] (17.04,-4.28) rectangle (17.76,-3.7199999999999998);
\draw[serieA!45, line width=0.6pt, rounded corners=0.6pt] (18.04,-4.28) rectangle (18.76,-3.7199999999999998);
\draw[serieA!45, line width=0.6pt, rounded corners=0.6pt] (19.04,-4.28) rectangle (19.76,-3.7199999999999998);
\fill[serieA, rounded corners=0.6pt] (20,-4.32) rectangle (20.8,-3.68);
\draw[serieA!45, line width=0.6pt, rounded corners=0.6pt] (21.04,-4.28) rectangle (21.76,-3.7199999999999998);
\draw[serieA!45, line width=0.6pt, rounded corners=0.6pt] (22.04,-4.28) rectangle (22.76,-3.7199999999999998);
\draw[serieA!45, line width=0.6pt, rounded corners=0.6pt] (23.04,-4.28) rectangle (23.76,-3.7199999999999998);
\fill[serieA, rounded corners=0.6pt] (24,-4.32) rectangle (24.8,-3.68);
\draw[serieA!45, line width=0.6pt, rounded corners=0.6pt] (25.04,-4.28) rectangle (25.76,-3.7199999999999998);
\draw[serieA!45, line width=0.6pt, rounded corners=0.6pt] (26.04,-4.28) rectangle (26.76,-3.7199999999999998);
\fill[serieA, rounded corners=0.6pt] (27,-4.32) rectangle (27.8,-3.68);
\draw[serieA!45, line width=0.6pt, rounded corners=0.6pt] (28.04,-4.28) rectangle (28.76,-3.7199999999999998);
\node[anchor=north west, font=\footnotesize\sffamily\color{tinta2}, align=left] at (0,-4.7) {lleno: sujeto de prueba en ese fold \quad contorno: sujeto de entrenamiento};
\fill[serieB] (0,-5.75) rectangle (0.8,-5.45); \node[anchor=west, font=\footnotesize\sffamily] at (1,-5.6) {Campaña 1};
\fill[serieC] (7,-5.75) rectangle (7.8,-5.45); \node[anchor=west, font=\footnotesize\sffamily] at (8,-5.6) {CFT70};
\fill[serieA] (14,-5.75) rectangle (14.8,-5.45); \node[anchor=west, font=\footnotesize\sffamily] at (15,-5.6) {Campaña 3};
\end{tikzpicture}

\caption[Folds de la validación por sujetos]{Folds de la validación por sujetos. Cada columna es un sujeto y cada fila
un fold; el cuadro lleno marca los sujetos de prueba de ese fold. Cada sujeto es de prueba exactamente una vez.}
\label{fig:folds}
\end{figure}

Dentro de cada fold hay dos particiones internas, también por sujetos. Para la parada temprana de las redes se apartan
4 sujetos de entrenamiento (y 7 personas de LifeSnaps para el modelo de Tierra). Para calibrar el ancho del rango, los
sujetos de entrenamiento se dividen en 4 grupos de validación cruzada.

\subsection{Intervalos por bootstrap de sujetos}

Los intervalos de confianza del 95\pct{} se calculan con un bootstrap de sujetos \parencite{efron1979,efron1993}: se
remuestrean personas, no días, porque los días de una persona no son independientes \parencite{field2007}. Cada
métrica es un cociente de sumas (errores, días fuera, escalones detectados), así que cada uno de los 1000 remuestreos
suma numerador y denominador de los sujetos elegidos y recién ahí divide. Las diferencias entre dos modelos usan los
mismos remuestreos para ambos, lo que las hace pareadas, y el intervalo es el de los percentiles 2,5 y 97,5.

\subsection{Cómo se evitó hacerse trampa}

\textcite{kapoor2023} mostraron que la fuga de información entre entrenamiento y prueba es una de las causas más
comunes de resultados que no se reproducen en el aprendizaje automático aplicado. Estas son las protecciones del
trabajo:
\begin{itemize}
\item Los sujetos de prueba no entran a ningún ajuste: ni al corrector, ni al ancho calibrado, ni al piso de la
  desviación, ni a la parada temprana, ni al ruido por variable que escala la pérdida.
\item El modelo de Tierra se entrena de nuevo en cada fold, sin el pre-reposo de los sujetos de prueba. El segundo run
  del kernel tenía esta fuga y el tercero la corrigió.
\item El ancho calibrado sale de residuos fuera de muestra, por validación cruzada dentro de los sujetos de
  entrenamiento.
\item Centro y ancho se separan con modelos cruzados, para no atribuir a la adaptación lo que es solo calibrar.
\item Los escalones tienen signo al azar y la detección se informa a favor y en contra de la deriva.
\item El ajuste personal de $k$ días solo ve los días anteriores al día $k \le 28$, y todo se evalúa desde el día 28.
\end{itemize}

\section{Infraestructura}

El código vive en el paquete \texttt{nova\_tl} de la rama \texttt{transfer-learning}, con 27 pruebas automáticas. Los
datos y el código se suben a Kaggle como un conjunto de datos privado, y el kernel corre las dos tareas en una GPU
\cifra{kernel.gpu}; el tercer run tardó \cifra{kernel.segundos}\un{s}. La regla del equipo es que las redes neuronales
se entrenan solo en Kaggle; el rango esperado, que es una regresión ridge, también se puede correr en un computador
personal. Los pesos se suben a Hugging Face desde el computador, porque la línea de comandos de Kaggle no maneja
secretos. Las librerías son PyTorch \parencite{paszke2019}, \texttt{transformers} \parencite{wolf2020},
scikit-learn \parencite{pedregosa2011}, NumPy \parencite{harris2020}, SciPy \parencite{virtanen2020} y pandas
\parencite{mckinney2010}, en las versiones de la \cref{tab:entorno}.

\begin{table}[htbp]
\caption[Entorno del tercer run]{Entorno del tercer run del kernel, tal como lo registró \texttt{metricas.json}.}
\label{tab:entorno}
\small
\begin{tabular}{ll}\toprule
 Componente & Versión \\
\midrule
pandas & \texttt{2.3.3} \\
numpy & \texttt{2.1.3} \\
scikit-learn & \texttt{1.6.1} \\
scipy & \texttt{1.16.3} \\
torch & \texttt{2.11.0+cu130} \\
transformers & \texttt{5.20.2} \\
gpu & \texttt{Tesla T4} \\
chronos-forecasting & \texttt{2.3.2} \\
\bottomrule\end{tabular}

\end{table}
